\chapter{Desenvolupament i implementació}
\label{cap:desarrollo}

Aquest capítol descriu l'\textbf{estat actual} de l'API i de l'app, no el prototip LangChain arxivat. El disseny del capítol~\ref{cap:arquitectura} es concreta aquí en etapes, rutes i algoritmes.

\section{Pipeline d'ingestió}

Ingerir un PDF és la fase de pujar del RAG híbrid (secció~\ref{sec:rag-dos-temps}): convertir un fitxer privat en fragments recuperables, amb estat visible i reintentable. El cas d'ús d'ingestió recorre sis etapes ordenades. Cada una persisteix una execució d'etapa (èxit, error, durada). Això permet el CLI de debug per rang i el botó \emph{Reintentar}, que torna a executar el pipeline sobre el mateix objecte a Storage.

\begin{figure}[htbp]
\centering
\begin{tikzpicture}[
  font=\small,
  box/.style={draw, rounded corners=2pt, align=center, inner sep=4pt, minimum height=0.85cm},
  arr/.style={-{Latex}, thick}
]
\node[box] (a) {download};
\node[box, right=0.35cm of a] (b) {parse};
\node[box, right=0.35cm of b] (c) {chunk};
\node[box, right=0.35cm of c] (d) {embed};
\node[box, right=0.35cm of d] (e) {store};
\node[box, right=0.35cm of e] (f) {synopsis};
\draw[arr] (a) -- (b);
\draw[arr] (b) -- (c);
\draw[arr] (c) -- (d);
\draw[arr] (d) -- (e);
\draw[arr] (e) -- (f);
\end{tikzpicture}

\caption{Etapes d'ingestió d'un PDF: del fitxer a l'índex (i, opcionalment, a la sinopsi).}
\label{fig:pipeline}
\end{figure}

\begin{description}
  \item[download] El worker llegeix els bytes des de Storage (\texttt{storage\_path} del document). Sense aquesta etapa, parsejaries un PDF que l'usuari ja no té o que no s'ha acabat de pujar.
  \item[parse] PyPDF extreu text \emph{per pàgina}. Un PDF escanejat (imatge) dóna cadenes buides: l'etapa falla de manera visible; no s'inventa OCR (fora d'abast).
  \item[chunk] El trossejador semàntic parteix el text preferint límits naturals: títols de capítol/secció, paràgrafs, oracions i, només al final, un tall dur de caràcters. Mida màxima 1200 caràcters, solapament 150. El solapament evita perdre una frase partida entre dos chunks. Cada fragment guarda pàgina inicial i final per a les cites.
  \item[embed] Cada chunk es converteix en un vector (OpenAI \texttt{text-embedding-3-small}, 1536 dimensions). Aquesta és la representació densa; el text cru es conserva per a FTS i per mostrar el snippet.
  \item[store] S'escriuen els chunks a \texttt{document\_chunks} (embedding + text). El document passa a estat llest.
  \item[synopsis] Si \texttt{ENABLE\_HIERARCHICAL\_RAG} és actiu, un LLM resumeix el document en 4--6 frases a \texttt{documents.synopsis}. Si el flag és fals, l'etapa es marca com a omesa, no com a error.
\end{description}

\subsection{Cua i worker}

Pujar un PDF no ha de bloquejar la petició HTTP uns minuts (parse + $N$ embeddings). El patró és una \textbf{cua de treballs}: \texttt{POST /documents/upload} desa el fitxer, crea un \texttt{Job} i torna de seguida; la UI consulta l'estat (\emph{en cua} / \emph{processant} / \emph{llest} / error).

Un \textbf{worker} és un bucle que reclama el següent job lliure. Aquí el worker d'ingestió corre al cicle de vida de FastAPI (engegada/aturada de l'app): cada segon reclama un job (migració 0007) i executa el pipeline. No hi ha procés worker a part al Compose: és més simple de desplegar en un TFM i més fràgil en producció (un restart mata la ingestió a mig fer; capítol~\ref{cap:conclusiones}).

Amb \texttt{ENABLE\_ASYNC\_INGESTION=true} (defecte) aquest bucle està actiu. Desactivar-lo serveix per depurar el pipeline a mà via CLI.

\begin{figure}[htbp]
\centering
\begin{tikzpicture}[
  font=\footnotesize,
  actor/.style={align=center, font=\small},
  arr/.style={-{Latex}, thick}
]
\node[actor] (ui) {UI};
\node[actor, right=2.1cm of ui] (api) {API};
\node[actor, right=2.1cm of api] (db) {Supabase};
\node[actor, right=2.1cm of db] (wk) {Worker};
\draw (ui.south) -- ++(0,-3.4);
\draw (api.south) -- ++(0,-3.4);
\draw (db.south) -- ++(0,-3.4);
\draw (wk.south) -- ++(0,-3.4);
\draw[arr] ([yshift=-0.45cm]ui.south) -- node[above, font=\tiny] {POST upload} ([yshift=-0.45cm]api.south);
\draw[arr] ([yshift=-0.95cm]api.south) -- node[above, font=\tiny] {PDF + job} ([yshift=-0.95cm]db.south);
\draw[arr] ([yshift=-1.55cm]wk.south) -- node[above, font=\tiny] {claim\_next} ([yshift=-1.55cm]db.south);
\draw[arr] ([yshift=-2.15cm]wk.south) -- node[above, font=\tiny] {pipeline 6 etapes} ([yshift=-2.15cm]db.south);
\draw[arr] ([yshift=-2.75cm]ui.south) -- node[above, font=\tiny] {GET job/doc} ([yshift=-2.75cm]api.south);
\end{tikzpicture}

\caption{Seqüència d'ingestió: la pujada torna de seguida; el worker reclama el job i escriu etapes.}
\label{fig:seq-ingest}
\end{figure}

\section{Xat, intents i cites}

El xat implementa la fase de preguntar del RAG híbrid (secció~\ref{sec:rag-pregunta}): autenticar, classificar, recuperar, construir un prompt que prohibeix inventar fora del context, generar, desar el fil. Hi ha dues rutes: JSON complet i streaming SSE (el servidor emet trossos de text a mesura que l'LLM els produeix, perquè la UI no esperi 10--20\,s en blanc).

La classificació d'intent (\texttt{classify\_query}) distingeix:

\begin{itemize}
  \item \emph{Resum}: ``de què va'', ``resumeix el temari''. Es reescriu la consulta i es diversifica per document perquè no guanyi un sol PDF.
  \item \emph{Capítol}: ``explica el tema 1'', ``capítulo 3''. Es filtra pels fitxers el nom dels quals comença pel número. Exemple: un PDF el nom del qual comença per \texttt{3\_} és el tema~3.
  \item \emph{Factual}: la resta; cerca híbrida estàndard. Només aquí l'agentic pot demanar una segona passada.
\end{itemize}

Les cites s'enriqueixen amb id de chunk, nom de fitxer, pàgines i un snippet de 280 caràcters. El frontend mostra el nom i obre un diàleg amb el fragment, renderitzat en Markdown. La política (prompts per intent, filtre administratiu, fusió amb chunks d'obertura) viu al mòdul de retrieval, no a la ruta HTTP.

El prompt factual (extracte) obliga l'ancoratge: \emph{``Answer using only the provided context. If the context is insufficient, say so. Cite sources as [n].''} Els intents de resum i capítol canvien l'èmfasi (cobrir tot el temari vs preferir un PDF numerat), no aquesta prohibició.

Un artefacte de flashcards es valida amb Pydantic abans de persistir; l'esquema mínim d'una carta és del tipus \texttt{\{"front": \ldots, "back": \ldots, "hint": \ldots\}} (pista opcional). Si l'LLM no respecta l'esquema, la petició falla: no es desa JSON arbitrari.

\begin{figure}[htbp]
\centering
\begin{tikzpicture}[
  font=\footnotesize,
  actor/.style={align=center, font=\small},
  arr/.style={-{Latex}, thick}
]
\node[actor] (ui) {UI};
\node[actor, right=2.0cm of ui] (api) {API / xat};
\node[actor, right=2.0cm of api] (ret) {Retrieval};
\node[actor, right=2.0cm of ret] (llm) {LLM};
\draw (ui.south) -- ++(0,-3.5);
\draw (api.south) -- ++(0,-3.5);
\draw (ret.south) -- ++(0,-3.5);
\draw (llm.south) -- ++(0,-3.5);
\draw[arr] ([yshift=-0.4cm]ui.south) -- node[above, font=\tiny] {POST ask/stream} ([yshift=-0.4cm]api.south);
\draw[arr] ([yshift=-1.0cm]api.south) -- node[above, font=\tiny] {intent + embed} ([yshift=-1.0cm]ret.south);
\draw[arr] ([yshift=-1.6cm]ret.south) -- node[above, font=\tiny] {dens + FTS + RRF} ([yshift=-1.6cm]api.south);
\draw[arr] ([yshift=-2.2cm]api.south) -- node[above, font=\tiny] {prompt + context} ([yshift=-2.2cm]llm.south);
\draw[arr] ([yshift=-2.8cm]llm.south) -- node[above, font=\tiny] {tokens SSE} ([yshift=-2.8cm]ui.south);
\end{tikzpicture}

\caption{Seqüència del xat en streaming: intent, recuperació híbrida, LLM i tokens SSE cap a la UI.}
\label{fig:seq-chat}
\end{figure}

\section{Generació i biblioteca d'estudi}

Generar materials no substitueix el xat: produeix artefactes \emph{persistents} a partir del mateix índex. \texttt{POST /api/generate/flashcards} i \texttt{/quiz} creen un \texttt{study\_artifact} amb títol automàtic, \texttt{origin=generated} i \texttt{content\_json} validat amb Pydantic (l'LLM no pot tornar un JSON arbitrari: si no passa l'esquema, la petició falla). Després s'aplica el filtre lecture-focused.

L'API de biblioteca (CRUD) permet que l'estudiant no depengui de ``l'última generació'':

\begin{itemize}
  \item Llistat per assignatura i tipus, ordenat per \texttt{updated\_at} (el conjunt tocat fa dos minuts puja).
  \item Creació manual, canvi de títol, esborrat.
  \item CRUD imbricat de cartes (anvers, revers, pista opcional) i de preguntes (quatre opcions, índex correcte, explicació).
\end{itemize}

Al frontend, \texttt{FlashcardsPanel} i \texttt{QuizPanel} comparteixen \texttt{ArtifactLibrary}: practicar o editar; importar/exportar JSON; etiquetes d'origen (IA / manual / importat). A la pràctica de cartes, si hi ha pista, es mostra amb \emph{Mostrar pista} o la tecla H.

\section{Planner i SM-2}

Un \textbf{sistema de repetició espaiada} (SRS) programa el pròxim repàs en funció de com de bé s'ha recordat l'ítem. SM-2 \parencite{wozniak1990supermemo} és l'algoritme clàssic de SuperMemo (Anki el va popularitzar): determinista, amb poc estat, auditable.

L'estat d'una carta (\texttt{SrsState}) guarda:

\begin{itemize}
  \item \texttt{ease} (facil·litat, inicial 2,5; mínim 1,3);
  \item \texttt{interval\_days} (dies fins al pròxim repàs);
  \item \texttt{repetitions} (encerts seguits);
  \item \texttt{next\_review\_at}.
\end{itemize}

\texttt{apply\_sm2} (domini pur) rep una qualitat $q \in [0,5]$. La UI mapeja quatre botons: \emph{de nou} $q=1$, \emph{difícil} $q=3$, \emph{bé} $q=4$, \emph{fàcil} $q=5$. Si $q<3$, es reinicien les repeticions i l'interval torna a 1 dia. Si $q\geq 3$: el primer encert espera 1 dia, el segon 6, i després $\mathrm{interval} \leftarrow \mathrm{round}(\mathrm{interval} \times \mathrm{ease})$; l'\emph{ease} s'actualitza amb la fórmula clàssica SM-2. El planner llista les cartes vençudes i persisteix cada ressenya.

\paragraph{Exemple numèric.}
Sobre una carta nova, \emph{bé} ($q=4$) deixa interval d'1~dia i una repetició. Un segon \emph{bé}: 6~dies. \emph{De nou} ($q=1$): zero repeticions i 1~dia. El test de domini comprova això; no és una mesura d'aprenentatge.

Posar SM-2 al domini (sense I/O) permet tests que comproven ``qualitat 1 $\Rightarrow$ interval 1'' sense Postgres.

\section{Perfil, flags i proveïdors}

El perfil persisteix el nom visible i el flag agentic per usuari. El mateix flag es pot forçar per entorn (defecte apagat): doble palanca (producte vs desplegament). Els endpoints de debug del pipeline només es registren si el flag de debug és actiu, per no exposar relançaments d'ingestió en un entorn compartit.

Proveïdors: OpenAI o Gemini per a l'LLM; embeddings OpenAI per defecte. El \emph{factory} d'adaptadors és l'únic lloc que instància l'SDK: el xat no canvia.

\section{Frontend: autenticació i dades}

Després del login, Supabase deixa cookies de sessió. Un component de sincronització copia el JWT cap a l'estat del client perquè el client HTTP l'enviï a l'API. TanStack Query invalida llistes d'assignatures, documents, jobs, fils i artefactes després de mutacions: la biblioteca no mostra un títol vell si s'acaba d'editar.

Les pàgines d'estudi són embolcalls prims que munten panells. \texttt{AppShell} ofereix la navegació: Documents, Chat, Flashcards, Quiz, Repàs.

\section{Qualitat i CI}

``Qualitat'' aquí no vol dir nota subjectiva del xat; vol dir que el codi es pot tornar a executar amb les mateixes comprovacions:

\begin{itemize}
  \item \textbf{ruff:} linter i format; evita errors trivials i inconsistència d'estil.
  \item \textbf{mypy --strict:} comprovació estàtica de tipus a 77 fitxers del backend. Un port mal implementat es detecta abans de runtime.
  \item \textbf{pytest} unitari: API amb fakes, retrieval, artefactes, SRS, chunker, rerank, ports, factory, entitats, worker, agregació de benchmark. No necessiten claus d'LLM.
  \item \textbf{Vitest:} utilitats pures del frontend (\texttt{citations}, \texttt{artifact-meta}).
  \item \textbf{Playwright smoke:} landing + login, sense credencials LLM (comprova que l'app arrenca, no el RAG).
  \item Integració real: \texttt{RUN\_REAL\_INTEGRATION=1} (assignatura, pujada, ingestió, xat). Opt-in, fora de la CI, perquè depèn de xarxa i de quota.
\end{itemize}

A GitHub Actions, dues pipelines (Python 3.12: ruff, mypy, pytest unitari; i \texttt{tsc}, Vitest, build i Playwright smoke) s'executen a cada \texttt{push}/\texttt{PR} sobre la branca principal quan canvia el camí corresponent. La CI no crida OpenAI.

\section{Desplegament local}

Docker Compose aixeca API (8000) i app (3000, \texttt{output: standalone} de Next.js: un servidor Node autosuficient, apte per a imatge Docker). Les migracions 0001--0008 s'apliquen al projecte Supabase, no dins del contenidor de l'API: l'esquema viu al git, l'instància Postgres és allotjada. El CLI de pipeline relança etapes; el de benchmark mesura retrieval. L'annex~\ref{anexo:despliegue} detalla el procediment.
